\section{Discussion}\label{sec:discussion}

\paragraph{Memory is architecture.}
The census supports the title claim mechanically. Parameter count alone explains
\kvcParRsq{} of the variance in $\log_{10}$ KV bytes per token on the deduplicated
set --- a weak predictor, as H2 anticipated --- while attention geometry and
mechanism indicators alone reach \kvcGeomRsq{}, and the combined model
\kvcFullRsq{} (Table~\ref{tab:mechanism}). The
dominant mechanisms are discrete, inspectable design knobs, not scaling
accidents: sharing KV heads across query heads
\cite{shazeer2019mqa,ainslie2023gqa}, bounding the cache with sliding windows
\cite{xiao2023streamingllm,jiang2023mistral}, compressing keys and values into
latents \cite{deepseekai2024v2}, and sharing one cache across layers
\cite{sun2024yoco,wu2024cla}. Because these are choices, equal-parameter models
diverge widely: the dense and MoE medians sit at \kvcMedKvppDense{} versus
\kvcMedKvppMoe{} bytes per token per $10^9$ parameters. A deployment that sizes
memory by parameter count is therefore wrong in both directions, and the error
lands on the constraint that binds first \cite{kwon2023paged}.

\paragraph{The 2025 bump is composition, mostly duplicated geometry.}
Median per-parameter KV cost under the pre-registered unit fell from
\kvcMedianKvppTwentyThree{} (2023) to \kvcMedianKvppTwentyFour{} (2024), rose to
\kvcMedianKvppTwentyFive{} in 2025, then fell to \kvcMedianKvppTwentySix{} in 2026
(Table~\ref{tab:overview}). Two composition effects, not architecture, drive the rise.
First, fine-tune series that reuse a 2023-era geometry with new names survive
family-level deduplication, and \kvcWorstGeoClusterRecent{} such rows from the single
most-duplicated geometry landed in 2025 or later; collapsing every geometry to one row
makes the series monotone (\kvcMedianKvppGeoTwentyThree{} to
\kvcMedianKvppGeoTwentyFour, \kvcMedianKvppGeoTwentyFive,
\kvcMedianKvppGeoTwentySix). Second, the remaining 2025 cohort skews small: absolute KV
per token fell (\kvcMedianKvTokTwentyFour{} to \kvcMedianKvTokTwentyFive) while
per-parameter KV rose, the signature of small checkpoints whose attention stacks stay
proportionally large --- and of multimodal checkpoints whose parameter counts include
vision towers. Across checkpoints the pooled trend stays negative under every counting
unit (family: $\rho=\kvcSpearmanRhoPP$, $p=\kvcSpearmanPPp$; geometry:
$\rho=\kvcSpearmanRhoGeo$). The planning lesson is that a yearly fleet median tracks
what publishers ship that year, deduplication unit included; in a year dominated by
small releases and re-released fine-tunes, per-parameter heuristics will under-provision
KV memory even while per-token architecture improves.

\paragraph{MoE dilutes KV per parameter, not KV per sequence.}
H3 was supported with the strongest effect in the census: MoE median
\kvcMedKvppMoe{} versus dense \kvcMedKvppDense{} bytes per token per $10^9$
parameters (rank-biserial \kvcMoeRB, $p=$\kvcMoeP). The mechanism is dilution:
experts add weights without adding attention layers, so the cache is set by the
shared trunk \cite{jiang2024mixtral,deepseekai2024v2}. Dilution is not
cheapness. Per-sequence KV bytes are unchanged by expert count, while total
weights crowd the memory pool: \kvcInfeasibleBigModels{} deduplicated models with known parameters exceed
128\,GiB at 4-bit weights under the modeled 6\,GiB runtime overhead and admit no KV at
all (\kvcWeightsAloneBig{} of them on weights alone; Table~\ref{tab:capacity}). On
unified-memory edge devices, where weights and KV share one pool, this asymmetry
decides what can be served at all \cite{rajesh2025edgemem}.

\paragraph{Why H4 failed.}
Absolute KV per token fell across the window --- the median dropped from
\kvcMedianKvTokTwentyThree{} (2023) to \kvcMedianKvTokTwentySix{} (2026) --- but
the KV-parity context is a ratio of weight bytes to KV bytes, so it moves only
when KV falls \emph{relative to} weights. Both moved at once. Window-heavy and
latent designs, and above all MoE weight growth, pushed parity out or removed
it: \kvcParityInfShareTwentySix\% of the 2026 fleet has bounded caches and never
reaches parity, and the share of models whose KV outgrows weights before 32k
context fell from \kvcDominantShareTwentyThree\% (2023) to
\kvcDominantShareTwentySix\% (2026). Against this, the 2025 small-model mass
pulled the yearly median back to \kvcParityMedTwentyFive{} tokens, below the
2023 median of \kvcParityMedTwentyThree{}. The pooled 2025--2026 median reached
only \kvcHFourRatio{} times the 2023 value under the pre-registered unit, far short of
the tenfold threshold, so H4 is \kvcHFourVerdict{} there --- but under geometry
deduplication the same statistic is \kvcHFourRatioGeo$\times$ and crosses the
threshold (Fig.~\ref{fig:parity}). The direction of the hypothesis was right; its
magnitude depends on the counting unit, because parity inherits weight-byte growth and
release mix, not KV savings alone.

\paragraph{Capacity levers.}
Under fixed budgets, geometry decides which lever pays. At 24\,GiB, small dense
checkpoints leave most of the pool for KV; at 128\,GiB the constraint flips to
weights, where \kvcInfeasibleBigModels{} models fit nothing even at 4-bit
(Table~\ref{tab:capacity}, Fig.~\ref{fig:capacity}). Weight quantization
\cite{frantar2022gptq,dettmers2022int8} and KV quantization
\cite{liu2024kivi,hooper2024kvquant} expand opposite sides of the same unified
pool, so the cheaper lever is a property of the specific attention geometry,
not of parameter count.
